\documentclass[10pt,a4paper]{amsart}
\usepackage{amsmath,amssymb,amsthm,graphicx}
% For dealing with references we use the comment environment
\usepackage{verbatim}
\newenvironment{reference}{\comment}{\endcomment}
%\newenvironment{reference}{}{}
\newenvironment{slogan}{\comment}{\endcomment}
\newenvironment{history}{\comment}{\endcomment}
% For commutative diagrams we use Xy-pic
\usepackage[all]{xy}

% We use 2cell for 2-commutative diagrams.
\xyoption{2cell}
\UseAllTwocells
\theoremstyle{plain}
\newtheorem{theorem}[subsection]{Theorem}
\newtheorem{proposition}[subsection]{Proposition}
\newtheorem{lemma}[subsection]{Lemma}

\theoremstyle{definition}
\newtheorem{definition}[subsection]{Definition}
\newtheorem{example}[subsection]{Example}
\newtheorem{exercise}[subsection]{Exercise}
\newtheorem{situation}[subsection]{Situation}

\theoremstyle{remark}
\newtheorem{remark}[subsection]{Remark}
\newtheorem{remarks}[subsection]{Remarks}

\numberwithin{equation}{subsection}
% Macros
%
\def\lim{\mathop{\mathrm{lim}}\nolimits}
\def\colim{\mathop{\mathrm{colim}}\nolimits}
\def\Spec{\mathop{\mathrm{Spec}}}
\def\Hom{\mathop{\mathrm{Hom}}\nolimits}
\def\Ext{\mathop{\mathrm{Ext}}\nolimits}
\def\SheafHom{\mathop{\mathcal{H}\!\mathit{om}}\nolimits}
\def\SheafExt{\mathop{\mathcal{E}\!\mathit{xt}}\nolimits}
\def\Sch{\mathit{Sch}}
\def\Mor{\mathop{\mathrm{Mor}}\nolimits}
\def\Ob{\mathop{\mathrm{Ob}}\nolimits}
\def\Sh{\mathop{\mathit{Sh}}\nolimits}
\def\NL{\mathop{N\!L}\nolimits}
\def\CH{\mathop{\mathrm{CH}}\nolimits}
\def\proetale{{pro\text{-}\acute{e}tale}}
\def\etale{{\acute{e}tale}}
\def\QCoh{\mathit{QCoh}}
\def\Ker{\mathop{\mathrm{Ker}}}
\def\Im{\mathop{\mathrm{Im}}}
\def\Coker{\mathop{\mathrm{Coker}}}
\def\Coim{\mathop{\mathrm{Coim}}}

% Boxtimes
%
\DeclareMathSymbol{\boxtimes}{\mathbin}{AMSa}{"02}

%
% Macros for moduli stacks/spaces
%
\def\QCohstack{\mathcal{QC}\!\mathit{oh}}
\def\Cohstack{\mathcal{C}\!\mathit{oh}}
\def\Spacesstack{\mathcal{S}\!\mathit{paces}}
\def\Quotfunctor{\mathrm{Quot}}
\def\Hilbfunctor{\mathrm{Hilb}}
\def\Curvesstack{\mathcal{C}\!\mathit{urves}}
\def\Polarizedstack{\mathcal{P}\!\mathit{olarized}}
\def\Complexesstack{\mathcal{C}\!\mathit{omplexes}}
% \Pic is the operator that assigns to X its picard group, usage \Pic(X)
% \Picardstack_{X/B} denotes the Picard stack of X over B
% \Picardfunctor_{X/B} denotes the Picard functor of X over B
\def\Pic{\mathop{\mathrm{Pic}}\nolimits}
\def\Picardstack{\mathcal{P}\!\mathit{ic}}
\def\Picardfunctor{\mathrm{Pic}}
\def\Deformationcategory{\mathcal{D}\!\mathit{ef}}
\usepackage{cleveref}
\crefname{originallemma}{original lemma}{original lemmas}
\Crefname{originallemma}{Original lemma}{Original lemmas}
\makeatletter
\newlabel{lemma-fully-faithful}{{0G0B}{source}{}{sourcepointer-1}{}}
\newlabel{lemma-fully-faithful@cref}{{[originallemma][0][]0G0B}{[1][0][]source}{}{}{}}
\makeatother
\crefname{originalproposition}{original proposition}{original propositions}
\Crefname{originalproposition}{Original proposition}{Original propositions}
\makeatletter
\newlabel{proposition-siblings-isomorphic}{{0FZW}{source}{}{sourcepointer-2}{}}
\newlabel{proposition-siblings-isomorphic@cref}{{[originalproposition][0][]0FZW}{[1][0][]source}{}{}{}}
\makeatother
\crefname{originalproposition}{original proposition}{original propositions}
\Crefname{originalproposition}{Original proposition}{Original propositions}
\makeatletter
\newlabel{properties-proposition-characterize-ample}{{01Q3}{source}{}{sourcepointer-3}{}}
\newlabel{properties-proposition-characterize-ample@cref}{{[originalproposition][0][]01Q3}{[1][0][]source}{}{}{}}
\makeatother
\crefname{originallemma}{original lemma}{original lemmas}
\Crefname{originallemma}{Original lemma}{Original lemmas}
\makeatletter
\newlabel{varieties-lemma-vanishin-h0-negative}{{0FD7}{source}{}{sourcepointer-4}{}}
\newlabel{varieties-lemma-vanishin-h0-negative@cref}{{[originallemma][0][]0FD7}{[1][0][]source}{}{}{}}
\makeatother
\title{Theorem fully faithful in equiv.tex}
\author{The Stacks Project Authors}
\date{stacks-project:a04446e57ec1fbc252a871afcec7752fb2807b14; GNU Free Documentation License 1.2 (as published in the Stacks Project COPYING)}
\begin{document}
\maketitle
\noindent\textit{Editorial note.} All mathematical statement and proof text is reproduced verbatim from the author's chapter equiv.tex of the Stacks Project, fixed at commit a04446e57ec1fbc252a871afcec7752fb2807b14. The pinned preamble is supplied as its own file; the wrapper class, the xr-hyper wiring and the \textbackslash{}externaldocument list are omitted because the assembler supplies the class.\par\medskip
\noindent\textit{Editorial note.} Numbering inside this standalone excerpt is editorial, not the full Stacks book print numbering.\par\medskip
\noindent\textit{Editorial note.} The candidate is from the corrected band-I follow-on screen (line C) of the extractor's stacks-candidates line: its author label is absent from the live master in every form searched and it belongs to none of the earlier judged pools. The H2 judgement and its reason were made in this module by reading the author's own statement and complete proof in the fixed checkout. Original candidate key: stacks:equiv.tex:theorem-fully-faithful.\par\medskip
\noindent\textit{Editorial note.} Reference labels in the proof resolve to navigation pointers to the original author source; these pointers do not reproduce proofs of the prerequisites.\par\medskip
\noindent\textit{Allowed standard prerequisites (editorial).}\begin{itemize}
\item Prerequisites named by the author inside the counted span. The counted span is reproduced verbatim, including the author's own citations. Its prerequisites are the earlier results of the same fixed Stacks Project version that the author cites there and in the chapter's lead-in.
\end{itemize}\medskip
\ifdefined\hypertarget\hypertarget{sourcepointer-1}{}\fi
\noindent\textit{Source reference.} Original lemma 0G0B: sources/equiv.tex, source lines 3247--3247. Navigation pointer used in the counted span; the span contains the pointer only and the cited result is not reproduced inside it. Registered as a non-load-bearing navigation stub in the established form for Stacks material. No bibliography citation is involved.\par\medskip
\ifdefined\hypertarget\hypertarget{sourcepointer-2}{}\fi
\noindent\textit{Source reference.} Original proposition 0FZW: sources/equiv.tex, source lines 2258--2258. Navigation pointer used in the counted span; the span contains the pointer only and the cited result is not reproduced inside it. Registered as a non-load-bearing navigation stub in the established form for Stacks material. No bibliography citation is involved.\par\medskip
\ifdefined\hypertarget\hypertarget{sourcepointer-3}{}\fi
\noindent\textit{Source reference.} Original proposition 01Q3: sources/properties.tex, source lines 4557--4557. Navigation pointer used in the counted span; the span contains the pointer only and the cited result is not reproduced inside it. Registered as a non-load-bearing navigation stub in the established form for Stacks material. No bibliography citation is involved.\par\medskip
\ifdefined\hypertarget\hypertarget{sourcepointer-4}{}\fi
\noindent\textit{Source reference.} Original lemma 0FD7: sources/varieties.tex, source lines 11075--11075. Navigation pointer used in the counted span; the span contains the pointer only and the cited result is not reproduced inside it. Registered as a non-load-bearing navigation stub in the established form for Stacks material. No bibliography citation is involved.\par\medskip
\begin{theorem}[Orlov]
\label{theorem-fully-faithful}
\begin{reference}
\cite[Theorem 2.2]{Orlov-K3}; this is shown in \cite{Noah}
without the assumption that $X$ be projective
\end{reference}
Let $k$ be a field. Let $X$ and $Y$ be smooth proper schemes over $k$
with $X$ projective over $k$. Any $k$-linear fully faithful exact 
functor $F : D_{perf}(\mathcal{O}_X) \to D_{perf}(\mathcal{O}_Y)$
is a Fourier-Mukai functor for some kernel in
$D_{perf}(\mathcal{O}_{X \times Y})$.
\end{theorem}

\begin{proof}
Let $F'$ be the Fourier-Mukai functor which is a sibling of $F$
as in Lemma \ref{lemma-fully-faithful}.
By Proposition \ref{proposition-siblings-isomorphic} we have $F \cong F'$
provided we can show that $\textit{Coh}(\mathcal{O}_X)$ has enough
negative objects. However, if $X = \Spec(k)$ for example, then
this isn't true. Thus we first decompose $X = \coprod X_i$
into its connected (and irreducible) components and we
argue that it suffices to prove the result for each of the
(fully faithful) composition functors
$$
F_i :
D_{perf}(\mathcal{O}_{X_i}) \to
D_{perf}(\mathcal{O}_X) \to
D_{perf}(\mathcal{O}_Y)
$$
Details omitted. Thus we may assume $X$ is irreducible.

\medskip\noindent
The case $\dim(X) = 0$. Here $X$ is the spectrum of a finite (separable)
extension $k'/k$ and hence $D_{perf}(\mathcal{O}_X)$
is equivalent to the category
of graded $k'$-vector spaces such that $\mathcal{O}_X$ corresponds to the
trivial $1$-dimensional vector space in degree $0$.
It is straightforward to see that any two
siblings $F, F' : D_{perf}(\mathcal{O}_X) \to D_{perf}(\mathcal{O}_Y)$
are isomorphic. Namely, we are given an isomorphism
$F(\mathcal{O}_X) \cong F'(\mathcal{O}_X)$
compatible the action of the $k$-algebra
$k' = \text{End}_{D_{perf}(\mathcal{O}_X)}(\mathcal{O}_X)$
which extends canonically to an isomorphism on any graded $k'$-vector space.

\medskip\noindent
The case $\dim(X) > 0$. Here $X$ is a projective smooth
variety of dimension $> 1$. Let $\mathcal{F}$ be a coherent
$\mathcal{O}_X$-module. We have to show there exists a
coherent module $\mathcal{N}$ such that
\begin{enumerate}
\item there is a surjection $\mathcal{N} \to \mathcal{F}$ and
\item $\Hom(\mathcal{F}, \mathcal{N}) = 0$.
\end{enumerate}
Choose an ample invertible $\mathcal{O}_X$-module $\mathcal{L}$.
We claim that $\mathcal{N} = (\mathcal{L}^{\otimes n})^{\oplus r}$
will work for $n \ll 0$ and $r$ large enough.
Condition (1) follows from
Properties, Proposition \ref{properties-proposition-characterize-ample}.
Finally, we have
$$
\Hom(\mathcal{F}, \mathcal{L}^{\otimes n}) =
H^0(X, \SheafHom(\mathcal{F}, \mathcal{L}^{\otimes n})) =
H^0(X, \SheafHom(\mathcal{F}, \mathcal{O}_X) \otimes \mathcal{L}^{\otimes n})
$$
Since the dual $\SheafHom(\mathcal{F}, \mathcal{O}_X)$ is torsion free, this
vanishes for $n \ll 0$ by Varieties, Lemma 
\ref{varieties-lemma-vanishin-h0-negative}. This finishes the proof.
\end{proof}

\end{document}
